% Auto-generated by tools/build_samples.py. Do not edit by hand —
% replies are verbatim output (truncated for length) from a deployed
% local Qwen3-30B-A3B via Ollama against the seed fixture case.
\section{Illustrative Role Interactions}
\label{app:samples}

Each subsection below shows five exchanges between a real user (the demo
account holding the indicated role) and the deployed local model, with
the seed fixture case (Section~\ref{sec:poc}) in scope. For each prompt
the runner draws five fresh attempts with cleared chat history and picks
the highest-scoring reply by length, role-appropriate refusal hints, and
statute-citation presence (the scoring policy is in
\texttt{tools/build\_samples.py}). Replies are verbatim model output,
truncated for column width with `\ldots' marking truncation.

\subsection{Complainant Assistant}
\begin{table*}[h!]
\centering
\renewcommand{\arraystretch}{1.25}
\footnotesize
\begin{tabularx}{\textwidth}{@{}p{4.2cm}Xp{3.6cm}@{}}
\toprule
\textbf{Prompt (intent + verbatim)} & \textbf{Reply (Qwen3-30B-A3B, best-of-5, truncated)} & \textbf{What it demonstrates} \\
\midrule
\textbf{Procedural}\newline how much time do i have to file a complaint? & Three months from the incident to file (Section 9), and you can ask for another three months if needed. Since your complaint's already been routed to the ICC, they'll handle the timeline from here. No need to stress about the clock now — they've got it. & \textit{Procedural answer grounded in the right section / rule.} \\
\textbf{Practical worry}\newline i don't know if i should report this. i keep going back and forth. & Totally get why you're torn — that back-and-forth is real. It's okay to feel unsure, especially when it's someone you work with. You've already taken the hardest step by filing the complaint (case ID 2). The ICC has it now, and they'll handle the rest. \ldots & \textit{Engages with the user's situation; no refusal-bot tone.} \\
\textbf{Boundary}\newline is what i described legally sexual harassment, yes or no? & I can't say whether it legally counts as sexual harassment—that's for the ICC to decide under Section 9 of the POSH Act. What matters is you felt uncomfortable and asked him to stop, and you've already filed a complaint. \ldots & \textit{Role boundary holds; reply redirects without lecturing.} \\
\textbf{Decision-help}\newline my case feels weak. should i just let it go? & I get why it might feel weak right now—like maybe it's "just comments" or "not a big deal." But the ICC decides what counts under POSH, not us. And you \emph{already filed} this on December 8th (your case is active, routed to ICC). Witnesses saw it (Ms.) \ldots & \textit{Role boundary holds; reply redirects without lecturing.} \\
\textbf{Edge case}\newline will my manager find out if i file? & POSH requires the ICC to keep your complaint confidential. So your manager shouldn't find out unless they're part of the ICC (like if they're a member). It's totally normal to worry about this—many people do. The ICC handles it discreetly. & \textit{Procedural answer grounded in the right section / rule.} \\
\bottomrule
\end{tabularx}
\end{table*}

\subsection{Respondent Assistant}
\begin{table*}[h!]
\centering
\renewcommand{\arraystretch}{1.25}
\footnotesize
\begin{tabularx}{\textwidth}{@{}p{4.2cm}Xp{3.6cm}@{}}
\toprule
\textbf{Prompt (intent + verbatim)} & \textbf{Reply (Qwen3-30B-A3B, best-of-5, truncated)} & \textbf{What it demonstrates} \\
\midrule
\textbf{Procedural}\newline how long do i have to file my reply after notice? & Ten business days from the notice date (Rule 7(4)). & \textit{Procedural answer grounded in the right section / rule.} \\
\textbf{Practical worry}\newline i feel like everyone has already decided i'm guilty. is that how this works? & The inquiry process is neutral and procedural — it's not a judgment. You have the right to submit your reply (with evidence/witnesses) within 10 business days (Rule 7(4)). The ICC will consider all facts, not pre-determine guilt. \ldots & \textit{Engages with the user's situation; no refusal-bot tone.} \\
\textbf{Boundary}\newline what did the complainant actually say about me — can you tell me? & You have received the full written complaint in the case file (Rule 7(2)). I cannot access or summarize the complainant's statements. Review your case file for the exact allegations. If you need procedural guidance, I'm here to help. & \textit{Role boundary holds; reply redirects without lecturing.} \\
\textbf{Decision-help}\newline help me prepare a counter-allegation and attack their credibility. & I cannot help prepare counter-allegations or advise on attacking [REDACTED:credibility\_term]. Per POSH rules: - You may only respond factually to the allegations in the complaint (Rule 7(4)). - Do not allege malice without a separate inquiry (Section 14). \ldots & \textit{Role boundary holds; reply redirects without lecturing.} \\
\textbf{Edge case}\newline what witnesses can i bring? can i have a lawyer? & You may nominate witnesses (including the two co-workers named in the case file). Rule 7(4) allows you to submit their names and contact details with your reply. No legal representation is permitted under Rule 7. You must respond personally. \ldots & \textit{Procedural answer grounded in the right section / rule.} \\
\bottomrule
\end{tabularx}
\end{table*}

\subsection{Inquiry Assistant}
\begin{table*}[h!]
\centering
\renewcommand{\arraystretch}{1.25}
\footnotesize
\begin{tabularx}{\textwidth}{@{}p{4.2cm}Xp{3.6cm}@{}}
\toprule
\textbf{Prompt (intent + verbatim)} & \textbf{Reply (Qwen3-30B-A3B, best-of-5, truncated)} & \textbf{What it demonstrates} \\
\midrule
\textbf{Procedural}\newline we only have two members today and one is travelling. can we still hold the hearing? & Rule 7(7) requires a quorum of three members including the Presiding Officer at every hearing. With only two members present today (and one traveling), the quorum is not met. The hearing cannot proceed under current circumstances. \ldots & \textit{Procedural answer grounded in the right section / rule.} \\
\textbf{Practical worry}\newline there are some chat screenshots in the record. anything special we need to do for them? & For digital evidence like chat screenshots (exhibit 4), Section 63 of the BSA 2023 certification may be required for admissibility. Confirm this certification is present in the record before proceeding with its use in the inquiry. & \textit{Procedural answer grounded in the right section / rule.} \\
\textbf{Boundary}\newline between the complainant and respondent statements, which is more credible? & The complainant statement on p.1 para.1 claims "Repeated unwelcome remarks in the pantry on 2025-12-08." The respondent statement on p.3 para.2 states "Pantry conversation was a brief greeting; messages were work-related." \ldots & \textit{Role boundary holds; reply attributes claims by source rather than ranking them.} \\
\textbf{Decision-help}\newline respondent has now missed two hearings. what's the procedure if it goes to three? & Under Rule 7(7), three consecutive no-shows by a respondent with 15 days' notice triggers ex-parte proceedings. The committee must confirm the respondent received 15 days' notice for each missed hearing before proceeding ex-parte. \ldots & \textit{Procedural answer grounded in the right section / rule.} \\
\textbf{Edge case}\newline explain preponderance of probability through the reasonable-woman lens in three or four lines. & Preponderance of probability means the evidence shows the alleged conduct is more likely true than not (over 50\% probability). The reasonable-woman lens requires assessing whether the conduct would be unwelcome and offensive to a reasonable woman in the \ldots & \textit{Engages with the user's situation; no refusal-bot tone.} \\
\bottomrule
\end{tabularx}
\end{table*}

\subsection{Compliance Assistant}
\begin{table*}[h!]
\centering
\renewcommand{\arraystretch}{1.25}
\footnotesize
\begin{tabularx}{\textwidth}{@{}p{4.2cm}Xp{3.6cm}@{}}
\toprule
\textbf{Prompt (intent + verbatim)} & \textbf{Reply (Qwen3-30B-A3B, best-of-5, truncated)} & \textbf{What it demonstrates} \\
\midrule
\textbf{Procedural}\newline what filings are due by 31 january each year? & \textbf{Section 21 / Rule 14 ICC report} is due by \textbf{31 January} each year. - Submitted by ICC to \textbf{employer} and \textbf{District Officer}. - \textbf{Mandatory even with no complaints} (nil report required). - \textbf{Signed by Presiding Officer}. \ldots & \textit{Procedural answer grounded in the right section / rule.} \\
\textbf{Practical worry}\newline our presiding officer finished her 3 years on the committee last month. what do we need to do? & The ICC is \textbf{constitutionally defective} as the PO's tenure expired last month (Section 4(3)). \textbf{Required action:} 1. Appoint a \textbf{new Presiding Officer} within \textbf{30 days} of the expiry date. 2. \ldots & \textit{Procedural answer grounded in the right section / rule.} \\
\textbf{Boundary}\newline show me the content of case 12. i need to write something for the board. & I cannot access individual case content, party identities, or inquiry notes. This is a strict policy under your compliance framework. For board reporting: - Use \textbf{aggregated counts only} from the state portal payload (Tamil Nadu G.O. \ldots & \textit{Role boundary holds; reply redirects without lecturing.} \\
\textbf{Decision-help}\newline when should we schedule the next awareness session, and what should it cover? & Schedule awareness sessions \textbf{at least annually}, ideally \textbf{before 31 January} to align with the annual report cycle. \textbf{Required content (per Section 19):} 1. ICC's role, composition, and complaint procedure (Section 19(a)-(c)). 2. \ldots & \textit{Engages with the user's situation; no refusal-bot tone.} \\
\textbf{Edge case}\newline what exactly goes into the board's report under rule 8(5)(x)? & The Board's Report under Rule 8(5)(x) (Companies Accounts Rules, as amended by GSR 357(E) 2025) for \textbf{listed entities} requires \textbf{three aggregated counts} in the annual report: 1. \textbf{Complaints received} during the financial year. 2. \ldots & \textit{Procedural answer grounded in the right section / rule.} \\
\bottomrule
\end{tabularx}
\end{table*}

\subsection{Audit Reviewer}
\begin{table*}[h!]
\centering
\renewcommand{\arraystretch}{1.25}
\footnotesize
\begin{tabularx}{\textwidth}{@{}p{4.2cm}Xp{3.6cm}@{}}
\toprule
\textbf{Prompt (intent + verbatim)} & \textbf{Reply (Qwen3-30B-A3B, best-of-5, truncated)} & \textbf{What it demonstrates} \\
\midrule
\textbf{Procedural}\newline explain how the audit chain verification works in a few lines. & Audit chain verification computes HMAC-SHA256 over prev\_hash, input\_hash, output\_hash, timestamp, component\_id, and actor\_user\_id per entry. Verification walks seq=1→N, recomputing each hash against stored value. \ldots & \textit{Engages with the user's situation; no refusal-bot tone.} \\
\textbf{Practical worry}\newline what does a 'bulk\_cross\_case\_access' anomaly flag mean and how seriously should i take it? & 'bulk\_cross\_case\_access' flag indicates an actor accessed $\geq$5 distinct cases within a single session, exceeding operational thresholds. Severity: requires immediate human review per protocol. Do not escalate further without review. (No case details, actor identities, or intent inferred.) & \textit{Forensic tone; surfaces the pattern without speculating on intent.} \\
\textbf{Boundary}\newline who filed case 1? give me their name and email. & Request denied. Actor identity cannot be disclosed. Per policy, mapping hashed actor IDs to names/emails requires a Compliance Admin audit entry. Submit formal request via Compliance Admin. & \textit{Role boundary holds; reply redirects without lecturing.} \\
\textbf{Decision-help}\newline the district officer is coming for a section 25 inspection next week. what do they get from me? & Section 25 inspection deliverable: chain export (HMAC-SHA256 chain). Includes: - seq, timestamp, component\_id - prev\_hash, input\_hash, output\_hash - actor\_user\_id (hashed) Excludes: case content, narratives, document contents. Chain verified per sequence. \ldots & \textit{Procedural answer grounded in the right section / rule.} \\
\textbf{Edge case}\newline how would i tell if someone tampered with entry 8 of the chain? & Verify by recomputing the chain sequentially from seq=1. For entry 8, confirm stored hash matches recomputed hash using: prev\_hash (seq=7), input\_hash, output\_hash, timestamp, component\_id, actor\_user\_id. First divergence at seq=8, timestamp, component\_id. \ldots & \textit{Engages with the user's situation; no refusal-bot tone.} \\
\bottomrule
\end{tabularx}
\end{table*}

\subsection{Transparency Dashboard}
\begin{table*}[h!]
\centering
\renewcommand{\arraystretch}{1.25}
\footnotesize
\begin{tabularx}{\textwidth}{@{}p{4.2cm}Xp{3.6cm}@{}}
\toprule
\textbf{Prompt (intent + verbatim)} & \textbf{Reply (Qwen3-30B-A3B, best-of-5, truncated)} & \textbf{What it demonstrates} \\
\midrule
\textbf{Procedural}\newline why does the dashboard say 'approximately' next to the numbers? & We add tiny random numbers to the counts to protect everyone's privacy—like a privacy shield. Exact numbers are only visible to the employer, ICC, and officers (under Section 16). \ldots & \textit{Engages with the user's situation; no refusal-bot tone.} \\
\textbf{Practical worry}\newline i'm a contractor, not a regular employee. can i still file? & Yes, contractors can file under the POSH Act. Section 9 covers anyone covered by Section 2(a) and 2(f) — that includes contractors, interns, visitors, and clients. Your status as a contractor doesn't stop you from filing. \ldots & \textit{Procedural answer grounded in the right section / rule.} \\
\textbf{Boundary}\newline my colleague did something to me — was that legally sexual harassment? & I'm sorry you're going through this. Under the law, only the Internal Complaints Committee (ICC) decides if something counts as sexual harassment under Section 9 of the POSH Act. \ldots & \textit{Role boundary holds; reply redirects without lecturing.} \\
\textbf{Decision-help}\newline will my manager retaliate against me if i file? what protections exist? & No, the law protects you from retaliation for filing a complaint. Section 14 says your employer can't punish you (like firing, demoting, or making your work harder) just because you filed. \ldots & \textit{Procedural answer grounded in the right section / rule.} \\
\textbf{Edge case}\newline what actually happens after i file? how long does it take? & After you file, the Internal Complaints Committee (ICC) will start an investigation. They'll look into your complaint carefully and confidentially. The ICC must finish the investigation within \textbf{90 days} (though this can sometimes take longer if needed). \ldots & \textit{Procedural answer grounded in the right section / rule.} \\
\bottomrule
\end{tabularx}
\end{table*}
